\documentclass[12pt]{article}
\usepackage[a4paper,margin=1in]{geometry}
\usepackage{amsmath,booktabs,setspace}
\usepackage{xcolor}
\usepackage{titlesec}

\titleformat{\section}{\large\bfseries}{}{0em}{}

\title{ECON 101: Macroeconomics \\ Quiz 4}
\author{Instructor: Muhebullah Karimzada}
\date{October 29,  2025}

\begin{document}
\maketitle

\section*{Multiple Choice (1 mark each)}

\begin{enumerate}
    \item \textbf{Which of the following best describes the rule of 70?}\\
    A. The number of years required for real GDP to double $\approx 70$ divided by the growth rate of real GDP per capita.\\
    B. Real GDP grows by 70 percent every year.\\
    C. The price level doubles every 70 years.\\
    D. Productivity rises by 70 percent when labour input doubles.\\[4pt]

    \item \textbf{Which of the following would \emph{most directly} raise labour productivity in the long run?}\\
    A. A temporary tax rebate to consumers.\\
    B. Increased spending on worker training and education.\\
    C. A higher nominal wage.\\
    D. An increase in the money supply.\\[4pt]

    \item \textbf{If the Canadian government implements policies that increase national saving, in the long run we would expect:}\\
    A. Lower investment and slower growth.\\
    B. Higher investment and faster growth.\\
    C. No change in potential output.\\
    D. Higher consumption and lower investment.\\[4pt]

    \item \textbf{If planned aggregate expenditure is less than actual real GDP, firms will …}\\
    A. Increase production.\\
    B. Decrease inventories.\\
    C. Decrease production.\\
    D. Raise prices to clear excess demand.\\[4pt]
\newpage
    \item \textbf{Which of the following would \emph{decrease} planned investment spending?}\\
    A. A decline in the real interest rate.\\
    B. A rise in expected future profitability.\\
    C. A deterioration in business confidence.\\
    D. An improvement in technology.\\[4pt]
\end{enumerate}


\noindent\textbf{Problem 1 (6 marks)}\\
Suppose a country’s real GDP per capita is \$50{,}000 and grows at 2 percent per year.

\begin{enumerate}
    \item[(a)] Using the rule of 70, approximate how many years it will take for real GDP per capita to double.\\
    \item[(b)] If productivity growth accelerates to 3 percent, how many years will doubling now take?\\
    \item[(c)] Explain in one or two sentences why even a small change in the growth rate matters for living standards over generations.
\end{enumerate}

\vspace{12pt}

\noindent\textbf{Problem 2 (9 marks)}\\
The following table shows a simplified economy (all figures in billions):

\begin{center}
\begin{tabular}{l l}
\toprule
\textbf{Component} & \textbf{Equation or Value} \\
\midrule
Consumption & $C = 250 + 0.8Y_d,\;\; Y_d = Y - T$\\
Planned Investment & $I = 150$\\
Government Purchases & $G = 200$\\
Net Exports & $NX = -50$\\
Taxes & $T = 100$\\
\bottomrule
\end{tabular}
\end{center}

\begin{enumerate}
    \item[(a)] Compute planned aggregate expenditure (AE) as a function of $Y$.\\
    \item[(b)] Find the equilibrium level of output where $AE = Y$.\\
    \item[(c)] At this equilibrium, determine household saving.\\
    \item[(d)] Briefly interpret whether the government is running a budget surplus or deficit at this output level.\\
\end{enumerate}

\end{document}
